\chapter{จลนศาสตร์และการเคลื่อนที่ของวัตถุ}
\label{ch:ch03_kinematics}
\index{จลนศาสตร์และการเคลื่อนที่}

\begin{conceptbox}[title=วัตถุประสงค์การเรียนรู้ประจำบท]
เมื่อสิ้นสุดการศึกษาบทเรียนนี้ นักศึกษาสามารถ
\begin{enumerate}
    \item อธิบายความหมายและความแตกต่างของปริมาณที่เกี่ยวข้องกับการเคลื่อนที่ ได้แก่ ตำแหน่ง การกระจัด ระยะทาง อัตราเร็ว ความเร็ว และความเร่งได้อย่างถูกต้อง
    \item เขียนและประยุกต์ใช้สมการการเคลื่อนที่แนวเส้นตรงด้วยความเร่งคงที่ (v = v₀ + at, s = v₀t + ½at$^2$, v$^2$ = v₀$^2$ + 2as) เพื่อแก้โจทย์ปัญหาได้
    \item วิเคราะห์การตกอย่างอิสระภายใต้สนามโน้มถ่วงโลกและคำนวณเวลา ความเร็ว หรือระยะทางที่เกี่ยวข้องได้
    \item วิเคราะห์การเคลื่อนที่แบบโพรเจกไทล์โดยแยกองค์ประกอบแนวราบและแนวดิ่งได้อย่างถูกต้อง และหาพิสัย เวลาที่ลอยอยู่ในอากาศ และความสูงสูงสุดได้
    \item อธิบายลักษณะของการเคลื่อนที่แบบวงกลมสม่ำเสมอ และคำนวณความเร่งสู่ศูนย์กลางได้
\end{enumerate}
\end{conceptbox}

\section{ทฤษฎีและหลักการพื้นฐาน}

\begin{bioappbox}
\textbf{จลนศาสตร์ของการตกตะกอนในเครื่องหมุนเหวี่ยงและอนุภาคในอากาศ}

ในห้องปฏิบัติการเทคโนโลยีชีวภาพ เครื่องหมุนเหวี่ยง (Centrifuge) ทำงานโดยอาศัยหลักการความเร่งสู่ศูนย์กลาง $a_c = v^2/r = \omega^2 r$ ซึ่งทำให้เกิดแรงเหวี่ยงหนีศูนย์กลางจำลองที่มีค่าสูงกว่าแรงโน้มถ่วงโลกหลายพันเท่า ($RCF \times g$) ส่งผลให้อนุภาคเซลล์ แบคทีเรีย หรือสารพันธุกรรม DNA ตกตะกอนแยกชั้นตามความหนาแน่นและมวลได้อย่างรวดเร็ว นอกจากนี้ สมการจลนศาสตร์แนวดิ่งยังใช้อธิบายอัตราการตกสัมพัทธ์ของสปอร์เชื้อราและละอองเกสรดอกไม้ในบรรยากาศ
\end{bioappbox}

\begin{labbox}
1. ตรวจสอบตำแหน่งศูนย์ (Zero Error) ของเครื่องมือวัดทุกชนิดก่อนเริ่มบันทึกข้อมูล
2. ทำการวัดซ้ำอย่างน้อย 3 ถึง 5 ครั้งในแต่ละจุดทดลองเพื่อคำนวณหาค่าเฉลี่ยและค่าเบี่ยงเบนมาตรฐาน
3. บันทึกผลการวัดพร้อมระบุหน่วยเอสไอ (SI Units) และค่าความคลาดเคลื่อนของการวัดเสมอ
\end{labbox}

\section{ตัวอย่างโจทย์และการวิเคราะห์เชิงฟิสิกส์}

\begin{examplebox}
ตัวอย่างที่ 3.1 รถยนต์คันหนึ่งเคลื่อนที่ด้วยความเร็วต้น 10 m/s ในแนวเส้นตรง แล้วเร่งเครื่องด้วยความเร่งคงที่ 2 m/s$^2$ เป็นเวลา 5 วินาที จงหาความเร็วสุดท้ายและระยะทางที่รถเคลื่อนที่ได้ในช่วงเวลานี้

วิธีทำ โจทย์กำหนด v₀ = 10 m/s, a = 2 m/s$^2$, t = 5 s

หาความเร็วสุดท้ายจาก v = v₀ + at v = 10 + (2)(5) = 10 + 10 = 20 m/s

หาระยะทางจาก s = v₀t + ½at$^2$ s = (10)(5) + ½(2)(5)$^2$ = 50 + (1)(25) = 50 + 25 = 75 m
\end{examplebox}

\section{แบบฝึกหัดทบทวนและคำถามท้ายบท}

\begin{enumerate}
    \item จงอธิบายความแตกต่างระหว่างการกระจัดกับระยะทาง พร้อมยกตัวอย่างสถานการณ์ที่ทั้งสองปริมาณมีค่าต่างกัน
    \item จงอธิบายความแตกต่างระหว่างอัตราเร็วเฉลี่ยกับความเร็วเฉลี่ย
    \item ความเร็วขณะหนึ่งหาได้จากกราฟตำแหน่ง–เวลาอย่างไร จงอธิบาย
    \item รถยนต์คันหนึ่งเคลื่อนที่ด้วยความเร็วต้น 5 m/s แล้วเร่งด้วยความเร่งคงที่ 3 m/s$^2$ เป็นเวลา 4 วินาที จงหาความเร็วสุดท้ายและระยะทางที่เคลื่อนที่ได้
    \item รถไฟเคลื่อนที่ด้วยความเร็ว 30 m/s แล้วเบรกด้วยความหน่วงคงที่จนหยุดนิ่งภายในระยะทาง 225 m จงหาความหน่วง (ความเร่งที่มีทิศตรงข้ามกับการเคลื่อนที่)
    \item วัตถุตกจากที่สูง 78.4 m โดยไม่มีความเร็วต้น จงหาเวลาที่ตกถึงพื้นและความเร็วขณะกระทบพื้น
\end{enumerate}

% สิ้นสุดบทเรียน
